% =====================================================================
%  Centralized macros. ALWAYS edit notation here, never inline.
% =====================================================================

% --------- packages ---------
\usepackage[utf8]{inputenc}
\usepackage[T1]{fontenc}
\usepackage{lmodern}
\usepackage{amsmath,amssymb,amsthm,mathtools}
\allowdisplaybreaks[1]
\usepackage{xcolor}
\usepackage{hyperref}
\hypersetup{colorlinks=true,linkcolor=blue!50!black,citecolor=blue!50!black,urlcolor=blue!50!black}
\usepackage[numbers,square,sort&compress]{natbib}
\usepackage{aliascnt}
\usepackage{cleveref}
\usepackage{microtype}
\usepackage{xspace}

% --------- theorem environments ---------
\theoremstyle{plain}
\newtheorem{theorem}{Theorem}[section]
\newaliascnt{proposition}{theorem}
\newtheorem{proposition}[proposition]{Proposition}
\aliascntresetthe{proposition}
\newaliascnt{lemma}{theorem}
\newtheorem{lemma}[lemma]{Lemma}
\aliascntresetthe{lemma}
\newaliascnt{corollary}{theorem}
\newtheorem{corollary}[corollary]{Corollary}
\aliascntresetthe{corollary}

\theoremstyle{definition}
\newaliascnt{definition}{theorem}
\newtheorem{definition}[definition]{Definition}
\aliascntresetthe{definition}
\newaliascnt{assumption}{theorem}
\newtheorem{assumption}[assumption]{Assumption}
\aliascntresetthe{assumption}
\newaliascnt{example}{theorem}
\newtheorem{example}[example]{Example}
\aliascntresetthe{example}

\theoremstyle{remark}
\newaliascnt{remark}{theorem}
\newtheorem{remark}[remark]{Remark}
\aliascntresetthe{remark}

% --------- spaces ---------
\newcommand{\R}{\mathbb{R}}
\newcommand{\N}{\mathbb{N}}
\newcommand{\Z}{\mathbb{Z}}
\newcommand{\E}{\mathbb{E}}
\newcommand{\PP}{\mathbb{P}}
\newcommand{\HH}{\mathcal{H}}        % main Hilbert space
\newcommand{\HN}{\mathcal{H}_N}      % truncated subspace
\newcommand{\X}{\mathcal{X}}
\newcommand{\Y}{\mathcal{Y}}

% --------- operators ---------
\DeclareMathOperator{\Tr}{Tr}
\DeclareMathOperator{\Var}{Var}
\DeclareMathOperator{\Cov}{Cov}
\DeclareMathOperator{\KL}{KL}
\DeclareMathOperator{\TV}{TV}
\DeclareMathOperator*{\argmin}{arg\,min}
\DeclareMathOperator*{\argmax}{arg\,max}

% --------- forward problem ---------
\newcommand{\G}{\mathcal{G}}         % forward operator
\newcommand{\noise}{\eta}
\newcommand{\sigobs}{\sigma}         % observational noise level
\newcommand{\udag}{{u^{\dagger}}}     % true parameter
\newcommand{\Gamobs}{\Gamma}          % noise covariance
\newcommand{\Phim}{\Phi}              % misfit (potential)

% --------- prior / learned drift ---------
% Wrap macros that contain super/subscripts in braces so users can attach
% further super/subscripts without triggering "Double subscript" errors.
\newcommand{\muz}{{\mu_{0}}}                % prior
\newcommand{\muy}{{\mu^{y}}}                % posterior given data y
\newcommand{\muyN}{{\mu^{y}_{N}}}           % discretized posterior
\newcommand{\drift}{{b_{\theta}}}            % Cameron--Martin drift
\newcommand{\driftN}{{b_{\theta,N}}}         % projected drift
\newcommand{\score}{{s_{\theta}}}            % score estimator (computational only)
\newcommand{\dH}{{d_{\mathrm{H}}}}    % Hellinger distance
\newcommand{\dTV}{{d_{\mathrm{TV}}}}  % total variation
\newcommand{\dW}{{W_{2}}}              % Wasserstein-2

% --------- norms and inner products ---------
\newcommand{\norm}[1]{\left\lVert#1\right\rVert}
\newcommand{\abs}[1]{\left\lvert#1\right\rvert}
\newcommand{\inner}[2]{\langle #1,\,#2\rangle}

% --------- shorthand for common phrases ---------
\newcommand{\eg}{e.g.,\xspace}
\newcommand{\ie}{i.e.,\xspace}
\newcommand{\cf}{cf.\xspace}

% --------- cleveref names ---------
\crefname{theorem}{Theorem}{Theorems}
\crefname{lemma}{Lemma}{Lemmas}
\crefname{proposition}{Proposition}{Propositions}
\crefname{corollary}{Corollary}{Corollaries}
\crefname{assumption}{Assumption}{Assumptions}
\crefname{definition}{Definition}{Definitions}
\crefname{remark}{Remark}{Remarks}
\crefname{equation}{Eq.}{Eqs.}
\crefname{section}{Section}{Sections}
\crefname{figure}{Figure}{Figures}
\crefname{table}{Table}{Tables}
